\section{Identifying strange hadrons with ALICE}

The ALICE experiment employs a topological reconstruction technique in the central rapidity region to identify weakly-decaying strange hadrons through their 
decay topologies into charged particles: $\Lambda \rightarrow \textnormal{p}\, \pi^{-}$, 
$\ks \rightarrow \pi^{+}\, \pi^{-}$, and cascade decay channel for $\Xi^{-} \rightarrow \Lambda\, \pi^{-} \rightarrow \textnormal{p}\, \pi^{-}\, \pi^{-}$,
along with their charge conjugates. The reconstructed decay particles are filtered through a set of kinematic and geometric selection criteria to ensure 
the daughter track characteristics comply with the expected decay topologies and to reduce combinatorial background. The strange hadrons of interest 
are subsequently identified through invariant-mass technique. For relatively short-lived resonance $\phi \rightarrow \textnormal{K}^{+}\, \textnormal{K}^{-}$ where 
topological selections are not required, signal extraction can be performed based solely on invariant-mass technique. 

The experimental setup consists of a central barrel covering $|\eta| < 0.9$ over the full azimuthal angle ($\varphi$), 
focusing on vertex reconstruction, tracking, and charged particle identification, accompanied by dedicated forward detectors.
Key subdetectors for identifying strange hadrons include the six-layered silicon-based Inner Tracking System (ITS) \cite{ITS:2014} and an inert 
gas filled Time Projection Chamber (TPC) \cite{TPC:2014}, which are used in particle identification, tracking, momentum measurement and particle 
trajectory evolution, enabling particle track reconstruction in 3D space. The Time-Of-Flight detector (TOF) \cite{TOF:2019}, in conjunction with the 
ITS, is utilized to reject out-of-bunch pileup. The V0 detector \cite{V0:2013} comprises of two sets of scintillator arrays on each side of the 
interaction point, spanning the pseudorapidity ranges $2.8 < |\eta| < 5.1$ (V0A) and $-3.7 < |\eta| < -1.7$ (V0C). By measuring the number of 
clusters in the two innermost layers of the ITS (Silicon Pixel Detector) and the signal amplitude in the V0's scintillator arrays, charged 
particle multiplicity can be estimated for mid- and forward pseudorapidities. A detailed description and performance of the ALICE detector can 
be found in Ref. \cite{AlicePerf:2014}.